% FAC-EN-U0064
% Source: fac_body.tex lines 2432--2463.
% Authority: PDF physical pp. 64--65; printed pp. 259--260
% (Chapter III, no. 66 complete).

\phantomsection\label{fac:III:3:no:66}
\textbf{66. Coherent algebraic sheaves on projective varieties.}

Let $V$ be a closed subvariety of the projective space
$X=\mathbf{P}_r(K)$, and let $\mathcal{F}$ be a coherent algebraic sheaf on
$V$. Extending $\mathcal{F}$ by zero outside $V$, we obtain a coherent
algebraic sheaf on $X$
(cf.\ \hyperref[fac:II:2:no:39]{no.~39}), denoted by $\mathcal{F}^X$; we
know that
$H^q(X,\mathcal{F}^X)=H^q(V,\mathcal{F})$. The results of
\hyperref[fac:III:3:no:65]{the preceding number} therefore apply to the
groups $H^q(V,\mathcal{F})$. We first obtain
(taking \hyperref[fac:III:1:no:52]{no.~52} into account):

\phantomsection\label{fac:III:3:no:66:theorem:1}
\textsc{Theorem 1.} \textit{The groups $H^q(V,\mathcal{F})$ are
finite-dimensional vector spaces over $K$, and are zero for
$q>\dim V$.}

In particular, for $q=0$, we have:

\phantomsection\label{fac:III:3:no:66:corollary}
\textsc{Corollary.} \textit{$\Gamma(V,\mathcal{F})$ is a finite-dimensional
vector space over $K$.}

(It is natural to conjecture that
\hyperref[fac:III:3:no:66:theorem:1]{the theorem above} holds for every
\textit{complete} variety, in the sense of Weil~[16].)

\phantomsection\label{fac:III:3:no:66:construction:twist}
Let $U'_i=U_i\cap V$; the $U'_i$ form an open covering
$\mathfrak{U}'$ of $V$. If $\mathcal{F}$ is an algebraic sheaf on $V$, let
$\mathcal{F}_i=\mathcal{F}(U'_i)$, and let $\theta_{ij}(n)$ be the
isomorphism from $\mathcal{F}_j(U'_i\cap U'_j)$ onto
$\mathcal{F}_i(U'_i\cap U'_j)$ defined by multiplication by
$(t_j/t_i)^n$. We denote by $\mathcal{F}(n)$ the sheaf obtained from the
$\mathcal{F}_i$ by gluing by means of the $\theta_{ij}(n)$. The operation
$\mathcal{F}(n)$ has the same properties as the one defined in
\hyperref[fac:III:2:no:54]{no.~54}, which it generalizes; in particular,
$\mathcal{F}(n)$ is canonically isomorphic to
$\mathcal{F}\otimes\mathcal{O}_V(n)$.

We have $\mathcal{F}^X(n)=\mathcal{F}(n)^X$. Applying
\hyperref[fac:III:2:no:55:theorem:1]{Theorem~1 of no.~55}, together with
\hyperref[fac:III:3:no:65:prop:7]{Proposition~7 of no.~65}, we obtain:

\phantomsection\label{fac:III:3:no:66:theorem:2}
\textsc{Theorem 2.} \textit{Let $\mathcal{F}$ be a coherent algebraic sheaf
on $V$. There exists an integer $m(\mathcal{F})$ such that, for every
$n\geq m(\mathcal{F})$:}

\phantomsection\label{fac:III:3:no:66:theorem:2:item:a}
\textit{(a) For every $x\in V$, the $\mathcal{O}_{x,V}$-module
$\mathcal{F}(n)_x$ is generated by the elements of
$\Gamma(V,\mathcal{F}(n))$;}

\phantomsection\label{fac:III:3:no:66:theorem:2:item:b}
\textit{(b) $H^q(V,\mathcal{F}(n))=0$ for every $q>0$.}

\phantomsection\label{fac:III:3:no:66:remark}
\textsc{Remark.} It is essential to observe that the sheaf
$\mathcal{F}(n)$ depends not only on $\mathcal{F}$ and $n$, but also on the
\textit{embedding} of $V$ into the projective space $X$. More precisely,
let $P$ be the principal fiber space $\pi^{-1}(V)$, whose structure group
is $K^*$; for an integer $n$, let $K^*$ act on $K$ by the formula
\phantomsection\label{fac:III:3:no:66:formula:action}
\[
(\lambda,\mu)\longrightarrow\lambda^{-n}\mu
\qquad\text{if}\qquad
\lambda\in K^*\quad\text{and}\quad\mu\in K.
\]

\phantomsection\label{fac:III:3:no:66:construction:bundle}
Let $E^n=P\times_{K^*}K$ be the fiber space associated with $P$, with
typical fiber $K$ and endowed with
\hyperref[fac:III:3:no:66:formula:action]{the preceding operators}; let
$\mathcal{S}(E^n)$ be the sheaf of germs of sections of $E^n$
(cf.\ \hyperref[fac:II:2:no:41]{no.~41}). Taking account of the fact that
the $t_i/t_j$ form a system of changes of charts for $P$, we immediately
verify that $\mathcal{S}(E^n)$ is canonically isomorphic to
$\mathcal{O}_V(n)$. The formula
\phantomsection\label{fac:III:3:no:66:formula:twist-bundle}
\[
\mathcal{F}(n)=\mathcal{F}\otimes\mathcal{O}_V(n)
=\mathcal{F}\otimes\mathcal{S}(E^n)
\]
then shows that the operation $\mathcal{F}\longmapsto\mathcal{F}(n)$
depends only on \textit{the class of the fiber space $P$ defined by the
embedding $V\to X$}. In particular, if $V$ is normal, $\mathcal{F}(n)$
depends only on the linear-equivalence class of the hyperplane sections of
$V$ in the embedding under consideration (cf.~[17]).
